% Part: first-order-logic
% Chapter: sequent-calculus
% Section: quantifier-rules

\documentclass[../../../include/open-logic-section]{subfiles}

\begin{document}

\olfileid{fol}{seq}{qrl}

\olsection{قواعد سورها}

\subsection{قواعد $\lforall$}

\begin{defish}
\Axiom$ !A(t), \Gamma \fCenter \Delta$
\RightLabel{\LeftR{\lforall}}
\UnaryInf$ \lforall[x][!A(x)],\Gamma \fCenter \Delta$
\DisplayProof
\hfill
\Axiom$ \Gamma \fCenter \Delta, !A(a) $
\RightLabel{\RightR{\lforall}}
\UnaryInf$ \Gamma \fCenter \Delta, \lforall[x][!A(x)]$
\DisplayProof
\end{defish}

در \LeftR{\lforall}، $t$ ترمی بسته است (یعنی متغیر ندارد). در
\RightR{\lforall}، $a$~!!a{constant}ی است که نباید در هیچ جای دنبالهٔ
پایینیِ قاعدهٔ \RightR{\lforall} رخ دهد. $a$ را \emph{متغیر ویژهٔ}
استنتاج \RightR{\forall} می‌نامیم.\footnote{گرچه $a$ در قاعدهٔ بالا
!!a{constant} است، اصطلاح «متغیر ویژه» را به کار می‌بریم. این کاربرد
دلایل تاریخی دارد.}

\subsection{قواعد $\lexists$}

\begin{defish}
\Axiom$ !A(a), \Gamma \fCenter \Delta $
\RightLabel{\LeftR{\lexists}}
\UnaryInf$ \lexists[x][!A(x)], \Gamma \fCenter \Delta$
\DisplayProof
\hfill
\Axiom$ \Gamma \fCenter \Delta, !A(t) $
\RightLabel{\RightR{\lexists}}
\UnaryInf$ \Gamma \fCenter \Delta, \lexists[x][!A(x)]$
\DisplayProof
\end{defish}

باز هم $t$~ترمی بسته است و $a$~!!a{constant}ی است که در دنبالهٔ
پایینیِ قاعدهٔ \LeftR{\lexists} رخ نمی‌دهد. $a$ را \emph{متغیر ویژهٔ}
استنتاج \LeftR{\lexists} می‌نامیم.

این شرط که متغیر ویژه در دنبالهٔ پایینیِ استنتاج
\RightR{\lforall} یا \LeftR{\lexists} رخ ندهد، \emph{شرط متغیر ویژه}
نامیده می‌شود.

\begin{explain}
قرارداد زیر را به یاد آورید: هرگاه $!A$ !!a{formula}ی باشد که
!!{variable}~$x$ در آن آزاد است، این امر را با نوشتن~$!A(x)$ نشان
می‌دهیم. در همان بافت، $!A(t)$ کوتاه‌شدهٔ~$\Subst{!A}{t}{x}$ است.
بنابراین می‌توانستیم قاعدهٔ \RightR{\lexists} را چنین نیز بنویسیم:
\begin{prooftree}
  \Axiom$\Gamma \fCenter \Delta, \Subst{!A}{t}{x}$
  \RightLabel{\RightR{\lexists}}
  \UnaryInf$\Gamma \fCenter \Delta, \lexists[x][!A]$
\end{prooftree}
توجه کنید که $t$ ممکن است از پیش در~$!A$ رخ دهد؛ برای نمونه، ممکن است
$!A$ همان~$\Atom{\Obj P}{t,x}$ باشد. پس استنتاج
$\Gamma \Sequent \Delta, \lexists[x][\Atom{\Obj P}{t,x}]$ از
~$\Gamma \Sequent \Delta, \Atom{\Obj P}{t,t}$ کاربرد درستی از
~\RightR{\lexists} است---می‌توان یک یا چند رخداد~$t$ در مقدمه را، و
نه لزوماً همهٔ آن‌ها را، با !!{variable} مقید~$x$ «جایگزین» کرد.
بااین‌حال، شرط‌های متغیر ویژه در \RightR{\lforall} و
~\LeftR{\lexists} ایجاب می‌کنند که !!{constant}~$a$ در~$!A$ رخ ندهد.
پس نمی‌توان $\Gamma \Sequent \Delta, \lforall[x][\Atom{\Obj P}{a,x}]$
را به‌درستی از $\Gamma \Sequent \Delta, \Atom{\Obj P}{a,a}$ با استفاده
از~$\RightR{\lforall}$ استنتاج کرد.
\end{explain}

\begin{explain}
در \RightR{\lexists} و \LeftR{\lforall} هیچ محدودیتی بر ترم~$t$ وجود
ندارد. از سوی دیگر، در قواعد \LeftR{\lexists} و
\RightR{\lforall}، شرط متغیر ویژه ایجاب می‌کند که !!{constant}~$a$
در دنبالهٔ بالایی در هیچ جایی بیرون از~$!A(a)$ رخ ندهد. این شرط برای
تضمین صحت دستگاه ضروری است؛ یعنی دستگاه فقط دنباله‌هایی را که معتبرند
!!{derive}s می‌کند. بدون این شرط، استنتاج زیر مجاز می‌بود:
\begin{prooftree}
  \Axiom$!A(a) \fCenter !A(a)$
  \RightLabel{*\LeftR{\lexists}}
  \UnaryInf$\lexists[x][!A(x)] \fCenter !A(a)$
  \RightLabel{\RightR{\lforall}}
  \UnaryInf$\lexists[x][!A(x)] \fCenter \lforall[x][!A(x)]$
  \DisplayProof\bottomAlignProof
  \qquad
  \Axiom$!A(a) \fCenter !A(a)$
  \RightLabel{*\RightR{\lforall}}
  \UnaryInf$!A(a) \fCenter \lforall[x][!A(x)]$
  \RightLabel{\LeftR{\lexists}}
  \UnaryInf$\lexists[x][!A(x)] \fCenter \lforall[x][!A(x)]$
\end{prooftree}
اما $\lexists[x][!A(x)] \Sequent \lforall[x][!A(x)]$ معتبر نیست.
\end{explain}

\end{document}
